\section{Positive product indices and a uniform Vieta estimate}
\label{sec:positive}

At a positive product index, every lower-state visit pays a factor of
at least $\lambda^{-1}$. This makes the leading discrepancy small.
Equality then places an actual integral Markov triple close to a
triple of commuting matrix traces. The latter follows the Euclidean
algorithm exactly. We follow it with actual Vieta mutations, using
the positive quotient formula to keep the error uniformly bounded.

\begin{figure}[tb]
 \centering
 % Vector figure source. Included directly by the manuscript; no external data.
\begin{tikzpicture}[
 box/.style={draw,rounded corners=2pt,align=center,text width=49mm,
              minimum height=13mm,inner sep=4pt,font=\small},
 route/.style={-{Stealth[length=2mm]},thick},
 control/.style={densely dashed,gray!75!black,{Stealth[length=1.5mm]}-{Stealth[length=1.5mm]}}
 ]
 \node[box] (actual) at (0,0) {Positive integral Markov triple\\$F=0$};
 \node[box] (terminal) at (7.6,0) {Positive integral terminal entry\\$L_0\geqslant1$};
 \node[box] (baseline) at (0,-2.7) {Commuting baseline triple\\$F=4$};
 \node[box] (zero) at (7.6,-2.7) {Zero exponent\\$B_0=2/3$};
 \draw[route] (actual) -- node[above,font=\footnotesize,align=center]
 {Vieta\\mutations} (terminal);
 \draw[route] (baseline) -- node[below,font=\footnotesize,align=center]
 {Euclidean\\subtraction} (zero);
 \draw[control] (actual) -- node[right,font=\footnotesize,align=left]
 {initial error forced\\small by trace equality} (baseline);
 \draw[control] (terminal) -- node[right,font=\footnotesize,align=left]
 {controlled error\\forces $L_0<1$} (zero);
\end{tikzpicture}

 \caption{The two trace evolutions in the proof. Here
 $F(X,Y,Z)=X^2+Y^2+Z^2-XYZ$ uses full traces, whereas $L_i$ and
 $B_i$ are traces divided by three. Both evolutions use Vieta mutation
 on their own invariant level. Error control compares them without
 identifying the levels; at the final step it contradicts positive
 integrality. The corresponding normalized forcing is $4/9$.}
 \label{fig:vieta-comparison}
\end{figure}

\begin{theorem}[The positive-index parameter bound]
\label{thm:positive-bound}
In the minimum-ray model \eqref{eq:gap-ray-data}--\eqref{eq:gap-positive-model},
let $h\geqslant2$, $1\leqslant t<h$, and let the gap bits be arbitrary
with sum $t$. If $l\geqslant1$ and $V_l=u_{\eps m}$ for an integer
$m\geqslant1$, then
\begin{equation}\label{eq:positive-final-bound}
 p<4^{h+1}.
\end{equation}
Both signs of the single index are included. The gap pattern need not
be mechanical.
\end{theorem}

\subsection{A small positive discrepancy determines rational exponents}

Use the quantities from Section~\ref{sec:gap}, and put
$$
 Q=3^{h-1}c^h,\qquad q=m-hl-t,\qquad
 \delta=\frac{a\lambda^q}{Q}-1,
 \qquad (a,b)=\begin{cases}(c,d),&\eps=1,\\(d,c),&\eps=-1.
 \end{cases}
$$
Lemma~\ref{lem:positive-discrepancy} gives $\delta>0$. The exact
equality also gives $\delta<\mathcal R_l/X_l$. A mixed path with $j$
lower positions has at least two transitions. Its relative weight
\eqref{eq:gap-relative-path} is at most $\sigma\lambda^{-j}$.
The all-lower path is bounded separately by $\lambda^{-h-2}$, because
$l,t\geqslant1$. Hence
\begin{equation}\label{eq:positive-discrepancy-bound}
 0<\delta<\frac4{9p^2}\{(1+\lambda^{-1})^h-1\}
                  +\lambda^{-h-2}.
\end{equation}
If $p\geqslant h$, then $h/\lambda<2/5$. The elementary geometric
majorant $\binom hj\leqslant h^j$ gives
$$
 (1+\lambda^{-1})^h-1<\frac{h/\lambda}{1-h/\lambda}
 <\frac{5h}{3\lambda}.
$$
Since $h\geqslant2$, we may bound the last term in
\eqref{eq:positive-discrepancy-bound} by
$\lambda^{-4}<16/(625p^4)$. Therefore
\begin{equation}\label{eq:positive-small-delta}
 0<\delta<\frac1{p^3}\left(\frac{8h}{27}+\frac{16}{625}\right)
 <\frac h{p^3}.
\end{equation}
For $p\geqslant3h$, define
\begin{equation}\label{eq:positive-kappa}
 e_0=\frac4D,\qquad \kappa=9K=1+e_0.
\end{equation}
Then $0<\delta<1/(3p^2)<e_0<1/8$.

The defining equation $a\lambda^q=(1+\delta)3^{h-1}c^h$ now fixes
the coefficients up to a small multiplicative error. If $a=c$, set
$n=h-1$; if $a=d$, set $n=h+1$. In either case
\begin{equation}\label{eq:positive-coefficient-normalization}
 \begin{gathered}
 c=\frac{\lambda^{q/n}}3\alpha,\qquad
 d=\frac{\lambda^{-q/n}}3\beta,\qquad \beta=\kappa/\alpha,\\
 \alpha=
 \begin{cases}
 (1+\delta)^{-1/n},&a=c,\\
 [\kappa/(1+\delta)]^{1/n},&a=d.
 \end{cases}
 \end{gathered}
\end{equation}
These are exact formulas, not asymptotic expansions. In the first case
$\alpha=(1+\delta)^{-1/n}$ satisfies $1-e_0<\alpha<1$, while
$\beta=\kappa(1+\delta)^{1/n}$ satisfies $1<\beta<(1+e_0)^2<1+3e_0$;
the upper bound for $\beta$ comes from
$(1+\delta)^{1/n}\leqslant1+\delta<1+e_0$, not from the lower bound on
$\alpha$, which would give only $(1+e_0)/(1-e_0)$.
In the second case $1<\alpha,\beta<1+e_0$, because $\delta<e_0$.
Thus, in both cases,
\begin{equation}\label{eq:positive-coefficient-errors}
 |\alpha-1|,|\beta-1|<3e_0,
 \qquad \max(\alpha/\beta,\beta/\alpha)<2.
\end{equation}
If $q=0$, these estimates would give
$u_0=(\alpha+\beta)/3<2(1+3e_0)/3<1$, contradicting integrality.

Put $s=\lambda^{1/n}$ and $j=|q|>0$. The minimum ratio in
\eqref{eq:gap-ray-data} gives
\begin{equation}\label{eq:positive-index-normalization}
 s^{2j-n}\leqslant
 \begin{cases}
 \beta/\alpha,&q>0,\\
 \alpha/\beta,&q<0.
 \end{cases}
\end{equation}
For $s\geqslant4$, the integer $2j-n$ cannot be positive.
Consequently $1\leqslant j\leqslant n/2$. In particular the case
$n=1$, which can occur when $h=2$, is already impossible. Every
remaining case has $n\geqslant2$.

Define the commuting baselines
\begin{equation}\label{eq:positive-baselines}
 B_i(s)=\frac{s^i+s^{-i}}3,\qquad i\geqslant0.
\end{equation}
We abbreviate these scalar quantities to $B_i$ when $s$ is fixed;
they are distinct from the generator matrix $B$. The actual neighboring
Markov triple is
$$
 x=u_0,\qquad z=u_{-\operatorname{sign}(q)},\qquad p=B_n.
$$
Its baseline indices are respectively $j,n-j,n$. For example, when
$q>0$, the two terms of $z$ have exponents $-(n-j)$ and $n-j$;
when $q<0$, the same statement holds with $\alpha,\beta$ interchanged.
Thus \eqref{eq:positive-coefficient-errors} gives, for both entries,
\begin{equation}\label{eq:positive-initial-errors}
 |L_i-B_i|<e_0(s^i+s^{-i}),\qquad i=j,n-j,
 \qquad L_n=p=B_n.
\end{equation}
The factor three in \eqref{eq:positive-coefficient-errors} cancels
the division by three in the Binet formulas.

\subsection{A positive quotient suppresses the mutation error}

\begin{lemma}[A uniform Vieta neighborhood]\label{lem:uniform-vieta}
Let $n\geqslant2$, $1\leqslant j<n$, and $s\geqslant4$.
There is no positive integral Markov triple with entries labeled by
the indices $(j,n-j,n)$ such that
\begin{equation}\label{eq:positive-vieta-hypothesis}
 |L_i-B_i(s)|\leqslant5s^{i-2n}
 \quad\hbox{for each of its three entries}.
\end{equation}
The baseline indices need not initially be in the order of the actual
integer entries.
\end{lemma}

\begin{proof}
For positive indices $a,b$, put $c=a+b$ and $d=|a-b|$.
The commuting baselines satisfy
\begin{equation}\label{eq:positive-baseline-identities}
 3B_aB_b-B_c=B_d,
 \qquad B_a^2+B_b^2-B_cB_d=\frac49.
\end{equation}
Both identities follow by expanding $s^i+s^{-i}$. In particular their
Markov defect is $4/9$, which must be retained in any error estimate.
For the actual triple, Vieta mutation instead gives
\begin{equation}\label{eq:positive-quotient-mutation}
 L_d=3L_aL_b-L_c=\frac{L_a^2+L_b^2}{L_c}.
\end{equation}
The first formula proves integrality; the second proves positivity.
The resulting triple is again Markov. No ordering of the actual
entries is needed for these conclusions.

At maximum baseline index $c=a+b$, we maintain
\begin{equation}\label{eq:positive-error-invariant}
 |L_i-B_i|\leqslant5s^{i-2c}
 \quad\hbox{for all three entries}.
\end{equation}
It holds initially at $c=n$. Exchange $a,b$ if necessary to have
$a\geqslant b$, so $d=a-b$. Write $E_i=L_i-B_i$. Combining
\eqref{eq:positive-baseline-identities} and
\eqref{eq:positive-quotient-mutation} gives the exact error equation
\begin{equation}\label{eq:positive-quotient-error}
 L_c(L_d-B_d)
 =2B_aE_a+E_a^2+2B_bE_b+E_b^2-B_dE_c+\frac49.
\end{equation}
The quadratic errors are included. Since $c\geqslant2$,
\begin{equation}\label{eq:positive-denominator}
 L_c\geqslant\left(\frac13-5s^{-2c}\right)s^c
 \geqslant\left(\frac13-\frac5{256}\right)s^c>0.
\end{equation}
For every $i\geqslant0$, including the possible terminal index zero,
$B_i\leqslant(2/3)s^i$. Equation~\eqref{eq:positive-error-invariant}
therefore bounds the numerator in \eqref{eq:positive-quotient-error} by
\begin{align*}
 &10s^{-2b}+\frac{20}{3}s^{-2a}
   +25s^{-2c}(s^{-2a}+s^{-2b})+\frac49\\
 &\qquad\leqslant\frac{50}{3}s^{-2}+50s^{-6}+\frac49\\
 &\qquad\leqslant\frac{50}{48}+\frac{50}{4096}+\frac49
 <5\left(\frac13-\frac5{256}\right).
\end{align*}
The last strict margin is $1303/18432>0$. Dividing by
\eqref{eq:positive-denominator} yields
\begin{equation}\label{eq:positive-new-error}
 |L_d-B_d|<5s^{-c}=5s^{d-2a}.
\end{equation}
If $a>b$, the new maximum baseline index is $a<c$. The new entry
satisfies exactly \eqref{eq:positive-error-invariant} at that index,
and the unchanged entries satisfy its weaker new bounds. Thus the
constant five does not grow with the number of mutations.

Repeated subtraction of positive indices terminates at $(g,g,2g)$,
where $g=\gcd(j,n)\geqslant1$. Make one further actual mutation.
The preceding estimates also hold for $a=b=g$ and $d=0$, giving
$$
 \left|L_0-\frac23\right|<5s^{-2g}\leqslant\frac5{16}.
$$
But $L_0$ is a positive integer and
$L_0<2/3+5/16=47/48<1$. This contradiction also covers the case
$j=n/2$, where the terminal mutation is the first move.
\end{proof}

\subsection{Completion of the parameter bound}

Suppose $p\geqslant4^{h+1}$. Then $p>3h$, so the coefficient
normalization above applies, and for either choice of $n\leqslant h+1$,
$$
 s=\lambda^{1/n}>4^{(h+1)/n}\geqslant4.
$$
If $n=1$ there is already a contradiction in
\eqref{eq:positive-index-normalization}. Otherwise $n\geqslant2$,
and the exact identity $\sqrt D=\lambda-\lambda^{-1}$ gives
$$
 e_0=\frac{4s^{-2n}}{(1-s^{-2n})^2}.
$$
For $i\geqslant1$, $s\geqslant4$ and $n\geqslant2$,
\begin{equation}\label{eq:positive-initial-uniform-bound}
 e_0(s^i+s^{-i})
 \leqslant\frac{17}{16}\frac4{(1-1/256)^2}s^{i-2n}
 =\frac{16384}{3825}s^{i-2n}<5s^{i-2n}.
\end{equation}
Thus \eqref{eq:positive-initial-errors} gives the prohibited triple
of Lemma~\ref{lem:uniform-vieta}. This proves
Theorem~\ref{thm:positive-bound}.

The role of integrality is concentrated at the end of the argument.
The commuting baselines follow the subtractive Euclidean algorithm to
$B_0=2/3$. Positive quotient mutation keeps an actual Markov orbit close
enough that its terminal entry would also lie below one. This controls
all positive product indices with a single error constant, rather than
estimating a growing sequence of differences of large traces.
